\documentclass[11pt]{article}
\usepackage[margin=1in]{geometry}
\usepackage[T1]{fontenc}
\usepackage{lmodern}
\usepackage{amsmath}
\usepackage{xcolor}
\usepackage{listings}
\usepackage{booktabs}
\usepackage{enumitem}
\usepackage{tcolorbox}
\usepackage[hidelinks]{hyperref}

\definecolor{codebg}{RGB}{247,247,249}
\definecolor{kw}{RGB}{0,90,160}
\definecolor{cmt}{RGB}{110,120,110}
\definecolor{str}{RGB}{160,60,60}

\lstset{
  language=Python,
  backgroundcolor=\color{codebg},
  basicstyle=\ttfamily\footnotesize,
  keywordstyle=\color{kw}\bfseries,
  commentstyle=\color{cmt}\itshape,
  stringstyle=\color{str},
  numbers=none, frame=single, framesep=5pt, rulecolor=\color{gray!40},
  breaklines=true, showstringspaces=false, columns=fullflexible,
  xleftmargin=6pt, xrightmargin=6pt, aboveskip=8pt, belowskip=8pt,
}

\newtcolorbox{keybox}[1]{colback=blue!4,colframe=blue!45,title={#1},
  fonttitle=\bfseries,boxrule=0.6pt,arc=2pt}

\title{\vspace{-1.5cm}\textbf{How We Mapped Every EEG Dataset to 22 Channels}\\[2pt]
\large A Student Guide to the DiffEEG Data Pipeline}
\author{EEGdiff\_V2 project}
\date{\today}

\begin{document}
\maketitle
\vspace{-0.6cm}

% ------------------------------------------------------------------
\section{The problem in one sentence}
Our model was born speaking \textbf{one exact ``language''}: every EEG clip it
reads must be \textbf{22 channels $\times$ 1280 numbers, at 256\,Hz (5 seconds)},
with the 22 electrodes always in the \textbf{same order}. Every public dataset
speaks a \emph{different} language --- different electrode count, names, order,
and sampling rate. \emph{Processing} means translating each dataset into the
model's one language.

% ------------------------------------------------------------------
\section{The 22-channel ``frame'' (the key idea)}
Picture a \textbf{muffin tin with 22 labeled cups}, in a fixed order:
\begin{center}\ttfamily\small
FP1 FP2 F3 F4 C3 C4 P3 P4 O1 O2 F7 F8 T3 T4 T5 T6 A1 A2 FZ CZ PZ ROC
\end{center}
The model learned ``cup \#13 is always T3.'' For any new dataset we must drop
each electrode into the \textbf{correct cup}. If an electrode is missing, we
leave that cup \textbf{empty (all zeros)}. If we put an electrode in the
\emph{wrong} cup, there is \textbf{no error message} --- the model just silently
produces bad results. That is why the \emph{order} matters as much as the names.

% ------------------------------------------------------------------
\section{The universal 4-step recipe}
Almost every dataset goes through the same four steps
(\texttt{dataset\_tools/prepare\_external\_eeg.py}).

\paragraph{Step A --- clean the electrode name} so that \texttt{"EEG T7-REF"}
and \texttt{"t7"} both become \texttt{"T7"}:
\begin{lstlisting}
def canonical_channel(name):
    cleaned = re.sub(r"[^A-Za-z0-9]", "", name.upper())  # drop spaces/dashes
    cleaned = cleaned.replace("EEG", "")                 # drop the "EEG" prefix
    return ALIASES.get(cleaned, cleaned)                 # apply the alias table
\end{lstlisting}

\paragraph{Step B --- the alias table} because different labs use different
names for the \emph{same} spot on the head:
\begin{lstlisting}
TARGET_CHANNELS = ["FP1","FP2","F3","F4","C3","C4","P3","P4","O1","O2","F7",
                   "F8","T3","T4","T5","T6","A1","A2","FZ","CZ","PZ","ROC"]
ALIASES = {
    "T7":"T3", "T8":"T4", "P7":"T5", "P8":"T6",   # modern name -> old name
    "TP9":"A1", "TP10":"A2", "FPZ":"FP1",
    "FPZCZ":"FZ", "PZOZ":"PZ",                     # sleep-EDF's bipolar pairs
}
\end{lstlisting}

\paragraph{Step C --- place into the 22 cups, zero-fill the rest.} This is the
actual ``mapping'':
\begin{lstlisting}
def project_channels(data, names):
    out = np.zeros((22, data.shape[1]), dtype=np.float32)   # start all-zero
    present = {}
    for i, name in enumerate(names):
        c = canonical_channel(name)
        if c in TARGET_CHANNELS and c not in present:
            present[c] = i                                  # this electrode is cup c
    for j, target in enumerate(TARGET_CHANNELS):
        if target in present:
            out[j] = data[present[target]]                  # copy signal into its cup
    return out, sorted(set(TARGET_CHANNELS) - set(present)) # + list of empty cups
\end{lstlisting}

\paragraph{Step D --- resample to 256\,Hz and cut into 5-second windows:}
\begin{lstlisting}
def resample_full(data, sfreq, target_sfreq):     # e.g. 512 Hz -> 256 Hz
    up, down = int(target_sfreq), int(round(sfreq))
    return resample_poly(data, up, down, axis=-1)
# then slide a 1280-sample window across the recording
\end{lstlisting}

That is the whole trick. Below is what was \textbf{special} for each dataset.

% ------------------------------------------------------------------
\section{The processing script for every dataset}

\subsection*{TUSZ \& TUAB --- the ``home'' datasets (exact-name match)}
These use the same clinical montage the model was trained on, so the names
already match (\texttt{EEG FP1-REF} $\rightarrow$ \texttt{FP1}) --- no aliases,
just extract the 22 by name, resample $250\!\rightarrow\!256$\,Hz, and window.
(Processed by a collaborator; the recipe is recorded in
\texttt{TUAB/processed\_tuab/tuab\_processing\_report.json}.)
\textbf{Result: 22/22 cups filled.}

\subsection*{Bonn --- only one electrode exists}
Bonn is a text file with a \textbf{single} signal (no electrode name), 173.61\,Hz.
We put it in cup \#1 (FP1) and leave the other 21 empty.
\begin{lstlisting}
def prepare_bonn(raw_root, writer, args):
    for txt in txt_files:
        label = "seizure" if stem.startswith("S") else "normal"
        data = np.loadtxt(txt)[None, :]                      # one channel
        data = resample_to_target(data, 173.61, 256, ...)    # -> 256 Hz
        framed = np.zeros((22, data.shape[1]))               # 22 empty cups
        framed[0] = data[0]                                  # put signal in FP1
        writer.add(split_name(stem), label, framed)
\end{lstlisting}
\textbf{Result: 1/22 cups filled.}

\subsection*{EEGMMIDB (motor imagery) --- 64 electrodes, keep 19}
A 64-electrode cap, 160\,Hz. The alias table maps 19 of them onto our frame;
the 3 clinical channels it lacks (A1, A2, ROC) stay empty. Labels come from the
EDF \emph{annotations} (T1 = left fist, T2 = right fist).
\begin{lstlisting}
def prepare_eegmmidb(raw_root, writer, args):
    for edf in ...:
        spans = [(onset, onset+dur, t1_or_t2_label) for each annotation]
        window_edf(edf, writer, "eegmmidb", subject, spans, "rest", args)
        # window_edf does: project_channels -> resample_full -> slide window,
        # tagging each window by which annotation its center falls in
\end{lstlisting}
\textbf{Result: 19/22 cups filled.}

\subsection*{TUEV (event types) --- same montage, but tiny event labels}
Channel mapping is the same \texttt{project\_channels}. The special part is
\emph{windowing}, and the confusing bit is worth stating clearly.

\begin{keybox}{The problem: it is the \emph{labels} that are 1 second, not the input}
The EEG \textbf{recordings} are long and normal (e.g.\ 10 minutes). What is short
is the \textbf{labels}: TUEV annotates the data in tiny \textbf{1-second,
per-channel} chunks. Each \texttt{.rec} line is
\texttt{channel, start\_sec, stop\_sec, label} --- e.g.\
\texttt{14, 1.4, 2.4, 5} means ``on channel 14, seconds 1.4--2.4, label 5
(artifact).'' So ``1 second'' is the \emph{granularity of the annotations}, not
the size of the input we feed the model (which is still 5\,s).
\end{keybox}

\noindent\textbf{Why that breaks the naive approach.} Our model eats
\textbf{5-second} windows.
\begin{itemize}[nosep,leftmargin=1.4em]
  \item \textbf{TUSZ / TUAB} label \emph{continuously} --- a whole recording (or
  a long span) is ``seizure''/``normal.'' So you chop it into a fixed grid
  (0--5\,s, 5--10\,s, \dots) and every window inherits a clear label. Easy.
  \item \textbf{TUEV} labels are \emph{sparse 1-second events} scattered through
  a 10-minute recording. Tile that same fixed 5-second grid and ask ``is there an
  event at this window's center?'' --- the little 1-second markers almost never
  line up with the grid cells, so you throw away $\sim$98\% of windows.
\end{itemize}

\noindent\textbf{The fix (event-driven windowing).} Instead of a fixed grid, take
each event and grab the \textbf{5 seconds centered on it}, labeling that window
with the event. Every labeled event then becomes exactly one usable 5-second
training window. (Because TUEV annotates \emph{per channel}, the same event
appears on several channels at once, so we first run \texttt{merge\_events} to
collapse those per-channel 1-second markers into single spans.)
\begin{lstlisting}
def extract_event_windows(events, total_sec, window_sec):
    spans = merge_events(events)                 # join touching same-label events
    for s0, s1, lab in spans:
        if (s1 - s0) <= window_sec:
            start = center the window on the event   # event-driven, not a grid
    ...
# then: projected, missing = project_channels(data, raw.ch_names)  # same 22 cups
\end{lstlisting}
\textbf{Result: 22/22 cups (with per-session exceptions).}

\subsection*{Sleep-EDFx --- only 2 bipolar channels}
The hardest case: only \textbf{2} EEG channels, both \emph{bipolar pairs}
(\texttt{Fpz-Cz}, \texttt{Pz-Oz}) that match nothing in a normal montage. The two
aliases \texttt{FPZCZ$\rightarrow$FZ} and \texttt{PZOZ$\rightarrow$PZ} rescue it
--- without them, \emph{all 22 cups end up empty} (which happened once and
produced garbage until we caught it). Sleep-stage labels come from a separate
Hypnogram file, paired by shared filename prefix.
\begin{lstlisting}
def prepare_sleep_edfx(raw_root, writer, args):
    for psg in psgs:
        hyp = next(psg.parent.glob(psg.name[:7] + "*Hypnogram.edf"), None)  # pair
        spans = [(onset, onset+dur, sleep_stage) for each annotation]
        window_edf(psg, writer, "sleep_edfx", subject, spans, "sleep", args)
        # project_channels maps Fpz-Cz -> FZ, Pz-Oz -> PZ, the rest stay zero
\end{lstlisting}
\textbf{Result: 2/22 cups filled.}

% ------------------------------------------------------------------
\section{Siena --- the special case (processed by Yehya)}
Siena is different from all the others: it was \textbf{not} run through our
pipeline. Our collaborator \textbf{Yehya} had already processed it, and his
output \emph{already matched our model's spec}: 22 channels, in the right order,
at 256\,Hz, in 1280-sample (5\,s) windows. So there was \textbf{no channel
mapping and no resampling to do}. Our script
\texttt{dataset\_tools/convert\_siena\_npz.py} only re-packages his \texttt{.npz}
files into our folder layout, with a clean \emph{patient-wise} train/test split.
\begin{lstlisting}
for patient, recs in by_patient.items():
    split = "test" if patient in test_patients else "train"   # no patient in both
    for f in recs:
        d = np.load(f); X, Y = d["signals"], d["labels"]
        for x, y in zip(X, Y):
            label = "seizure" if y == 1 else "normal"
            buffers[(split, label)].append(x)                 # x is already 22x1280
\end{lstlisting}

\begin{keybox}{Did we change the volts? (short answer: no)}
Yehya's Siena data is stored in \textbf{volts}, so the numbers are tiny
(standard deviation $\approx 2\times10^{-4}$). Our model was trained on a
different scale (standard deviation $\approx 230$, about a \textbf{million times
bigger}). \textbf{We did \emph{not} rescale or convert the volts} --- we kept
Yehya's raw values exactly as they were.

The fix happens at \textbf{normalization time}. Before the model sees a signal
we always z-score it: $\;x_{\text{norm}} = (x - \text{mean}) / \text{std}$. The
trick is \emph{which} mean and std we use:
\begin{itemize}[nosep,leftmargin=1.4em]
  \item \textbf{What we do (correct):} use \textbf{Siena's own} mean/std, which
  are also in volts ($\text{std}\approx 2\times10^{-4}$). Dividing volts by a
  volts-scale std cancels the unit and gives a clean signal with
  $\text{std}\approx 1$ --- exactly what the model expects.
  \item \textbf{What would break it (the bug we avoided):} using \emph{THUSZ's}
  microvolt-scale std ($\approx 230$) on Siena's volt data would give
  $2\times10^{-4}/230 \approx 10^{-6}$ --- an essentially flat, dead signal, and
  the model would predict garbage.
\end{itemize}
\textbf{In one line:} we did not touch the volts; we just normalized Siena with
\emph{its own} statistics, which makes the physical unit irrelevant. When running
any eval/finetune script on Siena, point \texttt{-{}-norm-stats-dir} at
\texttt{Siena/processed\_siena/normalization} (Yehya's stats), \emph{never} at
the THUSZ normalization folder.
\end{keybox}

\textbf{Result: 19/22 cups filled} (A1, A2, ROC are absent in Siena's montage).

% ------------------------------------------------------------------
\section{Summary}
\begin{center}
\small
\begin{tabular}{@{}llcl@{}}
\toprule
Dataset & Special handling & Cups filled & Script \\
\midrule
TUSZ / TUAB & exact-name match & 22/22 & collaborator's (see report) \\
Bonn        & single trace $\rightarrow$ FP1 & 1/22 & \texttt{prepare\_bonn} \\
EEGMMIDB    & alias 19; labels from annotations & 19/22 & \texttt{prepare\_eegmmidb} \\
Siena       & re-package; keep volts, use own stats & 19/22 & \texttt{convert\_siena\_npz.py} \\
TUEV        & event-centered windows & 22/22 & \texttt{prepare\_tuev.py} \\
Sleep-EDFx  & bipolar $\rightarrow$ midline aliases & 2/22 & \texttt{prepare\_sleep\_edfx} \\
\bottomrule
\end{tabular}
\end{center}

\noindent\textbf{The three golden rules to remember:}
\begin{enumerate}[nosep,leftmargin=1.5em]
  \item \textbf{Order matters as much as names.} A right-set-wrong-order montage
  fails silently.
  \item \textbf{Normalize with each dataset's own statistics.} That is what makes
  Siena's volts (or any scale) work without touching the raw numbers.
  \item \textbf{Split so the test set is genuinely unseen ---
  and so it matches whatever you compare against.} Patient-wise for clinical
  datasets (Siena, TUSZ); subject-wise for EEGMMIDB. A random window split leaks
  the same person --- and often the same recording --- into both sides and
  inflates the result.
\end{enumerate}

\end{document}
